\documentclass[aspectratio=169]{beamer}

% The 'handout' option collapses all overlays into single slides
%\documentclass[aspectratio=169, handout]{beamer}

\usepackage{amsmath, amssymb, amsthm, mathtools}
\usepackage{algorithm, algorithmic}
\usepackage{tikz}
\usetikzlibrary{arrows.meta,positioning,fit,calc,shapes}

\newtheorem{proposition}[theorem]{Proposition}

% --- Standard Beamer Theme ---
\usetheme{Madrid}
\usecolortheme{default}

\include{macros} 


\title[LinMDPs]{Linear MDPs}
\author[Sham \& Kiant\'{e}]{Kiant\'{e} Brantley \& Sham Kakade}
\date{}


\begin{document}

\begin{frame}
\titlepage  
\end{frame}

\begin{frame}[noframenumbering]
  \frametitle{Outline}
  \tableofcontents
\end{frame}

\section{Motivation \& Setting}

\begin{frame}\frametitle{Recap: From Bandits to MDPs}

\begin{itemize}
\item \textbf{Tabular MDPs (UCBVI):} We achieved $\tilde{O}(\textrm{poly}(H) \sqrt{S A K})$ regret. 
  \begin{itemize}
    \item Great, but scales poorly when the state space $\Scal$ is enormous.
  \end{itemize}
  
\vspace{0.4cm}
\pause
\item \textbf{Linear Bandits (LinUCB):} We achieved $\tilde{O}(d \sqrt{K})$ regret.
  \begin{itemize}
    \item No dependence on the number of arms! We exploit the feature dimension $d$.
  \end{itemize}

\vspace{0.4cm}
\pause
\item \textbf{Today's Goal:} Can we combine these ideas? 
  \begin{itemize}
    \item We want an algorithm for episodic MDPs with $\text{poly}(d, H, \sqrt{K})$ regret.
    \item \alert{Crucially:} Zero explicit dependence on $|\Scal|$ or $|\Acal|$.
  \end{itemize}
\end{itemize}

\end{frame}

\begin{frame}\frametitle{Episodic MDP Setting}

\begin{itemize}
\item \textbf{Horizon:} $H$ steps per episode, total $K$ episodes.
\item \textbf{State/Action:} Large state space $\Scal$, action space $\Acal$, fixed start $s_0$.
\item \textbf{Rewards:} Deterministic, known, and bounded: $r_h(s,a) \in [0,1]$.
  \begin{itemize}
    \item \textit{Note: Unknown rewards just require an extra Hoeffding bonus. We assume them known to isolate the difficulty of learning transitions.}
  \end{itemize}
\end{itemize}

\vspace{0.4cm}
\textbf{The Interaction Loop:}
In episode $k$, the learner chooses a policy $\pi^k = \{\pi_h^k\}_{h=0}^{H-1}$, and generates a trajectory:
\[
s_0^k = s_0, \qquad a_h^k = \pi_h^k(s_h^k), \qquad s_{h+1}^k \sim P_h^\star(\cdot \mid s_h^k, a_h^k)
\]

\vspace{0.2cm}
\textbf{Regret:}
\[
R_K \;:=\; \sum_{k=0}^{K-1}\bigl(V^\star - V^{\pi^k}\bigr)
\]

\end{frame}

\section{The Linear MDP Model}
\begin{frame}[noframenumbering]
  \frametitle{Outline}
  \tableofcontents[currentsection]
\end{frame}

% =======================
% Linear Transition Model
% =======================
\begin{frame}\frametitle{The Linear Transition Model}

We assume the transition kernel is \emph{linear in known features}.

%\vspace{0.25cm}
\begin{assumption}[Linear MDP]\label{assump:linear_mdp_slides}
There exists a known feature map $\phi:\Scal\times\Acal\to\R^d$ with
$\sup_{s,a}\|\phi(s,a)\|_2\le 1$, and unknown parameters
$\{\mu_h^\star\}_{h=0}^{H-1}$ with $\mu_h^\star\in\R^{|\Scal|\times d}$ such that for all
$h$ and $(s,a)$,
\[
P_h^\star(\cdot\mid s,a) \;=\; \mu_h^\star\,\phi(s,a).
\]
Equivalently, for each $s'\in\Scal$, \alert{
$P_h^\star(s'\mid s,a) \;=\; \mu_h^\star(s')^\top \phi(s,a)$}.
%where $\mu_h^\star(s')\in\R^d$ is the $s'$-th row of $\mu_h^\star$.
\end{assumption}

\vspace{0.15cm}
\pause
\textbf{Low-rank view (matrix form).}
View $P_h^\star$ as a matrix in $\R^{(SA)\times S}$ and
$\Phi\in\R^{(SA)\times d}$. Then
\[
P_h^\star \;=\; \Phi\,(\mu_h^\star)^\top,
\qquad\text{so}\qquad \mathrm{rank}(P_h^\star)\le d.
\]

\vspace{0.15cm}
\pause
\textbf{Normalization (used to bound coefficients).}
For all $f:\Scal\to\R$,
\[
\big\|(\mu_h^\star)^\top f\big\|_2 \;\le\; \sqrt{d}\,\|f\|_\infty.
\]

\end{frame}


% ==========================================
% Examples & intuition: what is a Linear MDP?
% ==========================================
\begin{frame}\frametitle{Examples \& Intuition: What is a Linear MDP?}

\small
\textit{Warm-up (tabular):} take $\phi(s,a)$ one-hot in $\R^{SA}$, so $d=SA$.

\vspace{0.25cm}
\pause
\textbf{Example 1: Latent mixture / ``topic'' model}
\begin{itemize}
  \item Known mixture weights: $\phi(s,a)\in\Delta(d)$.
  \item Unknown components: $\mu_{h,1}^\star,\dots,\mu_{h,d}^\star\in\Delta(\Scal)$.
  \item Mixture of $d$ base transition kernels:
  \[
  P_h^\star(\cdot\mid s,a)=\sum_{i=1}^d \phi_i(s,a)\,\mu_{h,i}^\star(\cdot).
  \]
\end{itemize}

\vspace{0.10cm}
\pause
\textbf{Example 2: Known state embedding + unknown linear latent dynamics}
\begin{itemize}
  \item Suppose we are given a \emph{state embedding} $\psi:\Scal\to\R^d$ (learned or engineered).
  \item Unknown dynamics act linearly in the embedding space: $W_h^\star\in\R^{d\times d}$,
  \[
  P_h^\star(s'\mid s,a)=\psi(s')^\top W_h^\star \phi(s,a)
  \qquad (\text{with normalization so this is a distribution}).
  \]
  \item Interpretation: next-state probabilities depend on $(s,a)$ only through $\phi(s,a)$,
  and the dependence on $s'$ is summarized by $\psi(s')$.
\end{itemize}

\end{frame}

% ======================================
% Magic property: linearizing one-step lookahead
% ======================================
\begin{frame}\frametitle{The ``Magic'' Property: Linear Backups}

\small
\textbf{Key consequence:} for any bounded $f:\Scal\to\R$, $(P_h^\star f)(s,a)$ is linear in $\phi(s,a)$.

\vspace{0.03cm}
\begin{lemma}[Linearization of $P_h^\star f$]\label{lem:linearize_Pf_slides}
Fix $h$ and bounded $f$. Let $w_{h,f}^\star := (\mu_h^\star)^\top f\in\R^d$. Then
\[
(P_h^\star f)(s,a)
:= \E_{s'\sim P_h^\star(\cdot\mid s,a)}[f(s')]
\;=\;
\phi(s,a)^\top w_{h,f}^\star.
\]
If $\|f\|_\infty\le H$, then $\|w_{h,f}^\star\|_2\le H\sqrt d$.
\end{lemma}

\vspace{0.03cm}
\pause
\textbf{Proof sketch:}
\[
(P_h^\star f)(s,a)=\sum_{s'} f(s')\,P_h^\star(s'|s,a)
=\sum_{s'} f(s')\,\mu_h^\star(s')^\top \phi(s,a),
\]
\[
=\Big(\sum_{s'} f(s')\,\mu_h^\star(s')\Big)^\top \phi(s,a)
=( (\mu_h^\star)^\top f)^\top \phi(s,a).
\]

\end{frame}

% ======================================
% Implication: strong completeness
% ======================================
\begin{frame}\frametitle{Implication: A Strong Form of Completeness}

\small
\textbf{Recall (Bellman completeness):}\;
$f\in\Fcal \Rightarrow \Tcal_h f \in \Fcal$ \quad (closure only for $f$ in the hypothesis class).

\vspace{0.15cm}
\pause
\textbf{Linear MDP is stronger (for the transition part):}\;
for each stage $h$,
\[
\forall\, f:\Scal\to\R \ \text{bounded},\qquad
(P_h^\star f)(s,a)=\E[f(s')\mid s,a]\in \mathrm{span}(\phi).
\]
Equivalently, there exists $w_{h,f}^\star\in\R^d$ such that
\[
(P_h^\star f)(s,a)=\phi(s,a)^\top w_{h,f}^\star \qquad \forall (s,a).
\]

\vspace{0.15cm}
\pause
\textbf{Why this matters for DP / least squares:}
\begin{itemize}
  \item $V_{h+1}^\star$ is nonlinear (max over actions), and $\widehat V_{h+1}^k$ is even messier (bonus + clipping).
  \item \alert{Still:} the \emph{regression target} $(s,a)\mapsto (P_h^\star \widehat V_{h+1}^k)(s,a)$ is always exactly linear in $\phi(s,a)$.
  \item So each stage reduces to a \emph{well-specified} linear regression problem for $w_{h,\widehat V_{h+1}^k}^\star$.
\end{itemize}

\end{frame}


% ======================================
% Planning with a known model
% ======================================
\begin{frame}\frametitle{Planning with a Known Model}

\small
Suppose we knew the transition parameters $\{\mu_h^\star\}_{h=0}^{H-1}$.

\vspace{0.15cm}
\textbf{Backward DP:}\quad $ V_H \equiv 0$, and for $h=H-1,\dots,0$,
\[
Q_h^\star(s,a)=r_h(s,a)+(P_h^\star V_{h+1}^\star)(s,a),
\qquad
V_h^\star(s)=\max_{a\in\Acal} Q_h^\star(s,a).
\]

\vspace{0.15cm}
\pause
\textbf{Linear MDP plug-in:}\quad for any $f$,
\[
(P_h^\star f)(s,a)=\phi(s,a)^\top(\mu_h^\star)^\top f.
\]
So the Bellman backup can be written as
\[
Q_h^\star(s,a)=r_h(s,a)+\phi(s,a)^\top w_h^\star,
\qquad
w_h^\star := (\mu_h^\star)^\top V_{h+1}^\star \in \R^d.
\]

\vspace{0.15cm}
\pause
\textbf{Key takeaway:} we never need the full $P_h^\star(\cdot\mid s,a)$.
To plan, we can estimate the \emph{linear functional}:
\[
(s,a)\mapsto (P_h^\star V_{h+1})(s,a)=\phi(s,a)^\top w_{h,V_{h+1}}^\star.
\]
Lin-UCBVI does this with ridge regression + an elliptical bonus.

\end{frame}

% =========================================================
% Stage-wise ridge regression (compact + sets up bonus cleanly)
% =========================================================
\begin{frame}\frametitle{Stage-wise Ridge Regression (Estimating $P_h^\star f$)}

\small
At episode $k$, stage $h$, we have past transitions
$D_h^k=\{(s_h^i,a_h^i,s_{h+1}^i)\}_{i<k}$ and features $x_h^i:=\phi(s_h^i,a_h^i)\in\R^d$.

\vspace{0.10cm}
\textbf{Design matrix:}\;
\[
\Lambda_h^k := \lambda I + \sum_{i=0}^{k-1} x_h^i(x_h^i)^\top.
\]

\vspace{0.10cm}
\textbf{Ridge fit for a given target function $f:\Scal\to\R$:}\;
\[
\widehat w_{h,f}^k := (\Lambda_h^k)^{-1}\sum_{i=0}^{k-1} x_h^i\, f(s_{h+1}^i),
\qquad
(\widehat P_h^k f)(s,a):=\phi(s,a)^\top \widehat w_{h,f}^k.
\]

\vspace{0.10cm}
\pause
\textbf{Interpretation:}
we never estimate $P_h^\star(\cdot\mid s,a)$; we only estimate the \emph{linear functional}
$(s,a)\mapsto (P_h^\star f)(s,a)$ for the specific $f$ used in planning (namely $f=\widehat V_{h+1}^k$).

\end{frame}


% =========================================================
% Bonus: where it comes from (replaces the "Ridge and Bonus?" placeholder)
% =========================================================
\begin{frame}\frametitle{UCB Bonus: Uncertainty in Direction $\phi(s,a)$}

\small
Under the linear MDP model, for each fixed $f$ there is a true coefficient
$w_{h,f}^\star=(\mu_h^\star)^\top f$ with $(P_h^\star f)(s,a)=\phi(s,a)^\top w_{h,f}^\star$.
Ridge gives $\widehat w_{h,f}^k$.

\vspace{0.10cm}
\pause
\textbf{High-probability form (LinUCB-style):}\;
for all $(s,a)$,
\[
\big|\phi(s,a)^\top(\widehat w_{h,f}^k-w_{h,f}^\star)\big|
\ \le\
\beta\;\|\phi(s,a)\|_{(\Lambda_h^k)^{-1}}.
\]

\vspace{0.10cm}
\pause
\textbf{So we use the exploration bonus:}\;
\[
b_h^k(s,a) := \beta\,\|\phi(s,a)\|_{(\Lambda_h^k)^{-1}}.
\]
This is large when $\phi(s,a)$ lies in a poorly-sampled direction (small eigenvalues of $\Lambda_h^k$),
and shrinks as we collect more data.

\end{frame}


% =========================================================
% Optimistic planning (minor compacting)
% =========================================================
\begin{frame}\frametitle{Optimistic Planning (Lin-UCBVI)}

\small
\textbf{Exploration bonus:}\;
\[
b_h^k(s,a) := \beta\,\|\phi(s,a)\|_{(\Lambda_h^k)^{-1}}.
\]

\vspace{0.10cm}
\pause
\textbf{Backward induction:}\; set $\widehat V_H^k\equiv 0$. For $h=H-1,\dots,0$,
\begin{align*}
\onslide<2->{ \widehat Q_h^k(s,a) &:= \mathrm{clip}_{[0,H-h]}\!\Big( r_h(s,a) + b_h^k(s,a) + (\widehat P_h^k \widehat V_{h+1}^k)(s,a) \Big) \\ }
\onslide<3->{ \widehat V_h^k(s) &:= \max_{a\in\Acal}\widehat Q_h^k(s,a) \\ }
\onslide<4->{ \pi_h^k(s) &\in \argmax_{a\in\Acal}\widehat Q_h^k(s,a). }
\end{align*}

\onslide<5->{
\vspace{0.10cm}
Execute $\pi^k$ for one episode, then update $\Lambda_h^{k+1}=\Lambda_h^k+\phi(s_h^k,a_h^k)\phi(s_h^k,a_h^k)^\top$.
}

\end{frame}


\begin{frame}\frametitle{The Regret Guarantee}

\small
\begin{theorem}[Regret of Lin-UCBVI]
Assume the Linear MDP model and $r_h(s,a)\in[0,1]$. Run Lin-UCBVI with $\lambda\ge 1$ and
a bonus scale $\beta=\widetilde O(Hd)$ (from a uniform confidence bound).

Then with probability at least $1-\delta$, the cumulative regret satisfies:
\[
R_K \;=\; \sum_{k=0}^{K-1} \left(V^\star - V^{\pi^k}\right)
\;\le\;
c\,\beta\,H\,\sqrt{K\,d\,\log\!\Big(1+\tfrac{K}{\lambda}\Big)}
\;+\;
\mathcal{O}\Big(H\sqrt{K\log(1/\delta)}\Big),
\]
for a universal constant $c$.
\end{theorem}

\vspace{0.15cm}
\pause
\textbf{Takeaways:}
\begin{itemize}
  \item The $\mathcal{O}(\sqrt{K})$ martingale noise term is dominated by the primary $\beta$ term.
  \item Absorbing lower-order terms gives: $R_K \;=\; \widetilde O\!\Big(H^2\sqrt{d^3K}\Big)$.
  \item No explicit dependence on $|\Scal|$ or $|\Acal|$ (structure captured entirely by $d$).
\end{itemize}

\end{frame}


% =========================================================
% Top-level proof roadmap (new)
% =========================================================
\begin{frame}\frametitle{Proof Roadmap (What are the three moving parts?)}

\small
We reuse the same three ideas as UCBVI + LinUCB:

\vspace{0.10cm}
\begin{enumerate}
\item \textbf{(Confidence)} \quad
Uniformly control value-dependent regression targets:
\[
\big|(\widehat P_h^k f)(s,a)-(P_h^\star f)(s,a)\big|
\ \le\ \beta\,\|\phi(s,a)\|_{(\Lambda_h^k)^{-1}}
\quad \text{for } f=\widehat V_{h+1}^k.
\]
\item \textbf{(Optimism)} \quad
Backward induction shows $\widehat V_h^k(s)\ge V_h^\star(s)$ for all $k,h,s$.
\item \textbf{(Simulation + summation)} \quad
For each episode $k$,
\[
V^\star - V^{\pi^k}\ \le\ 2\sum_{h=0}^{H-1}\E_{d_h^{\pi^k}}\!\big[b_h^k(s_h,a_h)\big],
\]
then sum using an \textbf{elliptical potential} (log-det) argument.
\end{enumerate}

\vspace{0.10cm}
\pause
\textbf{Today’s only new wrinkle vs LinUCB:} the target $f=\widehat V_{h+1}^k$ is random, so we need
\emph{uniform} confidence over the whole optimistic value-function class.

\end{frame}


% =========================================================
% Statistical hurdle (tighten wording)
% =========================================================
\begin{frame}\frametitle{The Statistical Hurdle: Why Uniform Confidence?}

\small
In LinUCB we estimate a fixed parameter $\mu^\star$. Here we estimate a \emph{value-dependent} coefficient:
\[
\widehat w_{h,\widehat V_{h+1}^k}^k
=
(\Lambda_h^k)^{-1}\sum_{i=0}^{k-1}\phi(s_h^i,a_h^i)\,\widehat V_{h+1}^k(s_{h+1}^i).
\]

\vspace{0.12cm}
\pause
\textbf{The catch:}
\begin{itemize}
  \item $\widehat V_{h+1}^k$ is a \alert{random function} built from past data (and bonuses).
  \item So the ``labels'' $\widehat V_{h+1}^k(s_{h+1}^i)$ are correlated with the design/history.
  \item A fixed-function concentration bound is not enough; we need a bound that holds
  \emph{simultaneously for all optimistic value functions the algorithm may produce}.
\end{itemize}

\end{frame}


% =========================================================
% Uniform confidence (cleaner + less misleading note)
% =========================================================
\begin{frame}\frametitle{Step (a): Uniform Confidence over Optimistic Values}

\small
\begin{proposition}[Uniform model confidence (informal)]
With probability at least $1-\delta$, simultaneously for all episodes $k$, stages $h$,
all $(s,a)$, and the specific $f=\widehat V_{h+1}^k$ produced by the algorithm,
\[
\big|(\widehat P_h^k f)(s,a) - (P_h^\star f)(s,a)\big|
\ \le\ 
\beta\,\|\phi(s,a)\|_{(\Lambda_h^k)^{-1}}.
\]
\end{proposition}

\vspace{0.10cm}
\pause
\textbf{Where does $\beta$ come from?}
Self-normalized martingale concentration + an $\epsilon$-net over the optimistic value class
(max + bonus + clipping), which yields $\beta=\widetilde O(Hd)$.

\end{frame}


% =========================================================
% Step (b): Optimism (short)
% =========================================================
\begin{frame}\frametitle{Step (b): Optimism (Backward Induction)}

\small
On the uniform confidence event, for every episode $k$, stage $h$, and state $s$,
\[
\widehat V_h^k(s)\ \ge\ V_h^\star(s).
\]

\vspace{0.12cm}
\pause
\textbf{Reason (one line):}
for any $(s,a)$,
\[
(P_h^\star \widehat V_{h+1}^k)(s,a)
\le
(\widehat P_h^k \widehat V_{h+1}^k)(s,a)+b_h^k(s,a),
\]
so $r_h + P_h^\star V_{h+1}^\star \le r_h + b_h^k + \widehat P_h^k\widehat V_{h+1}^k$ (and clipping preserves the inequality).

\end{frame}


% =========================================================
% Step (c): Simulation lemma + visitation measure (make it explicit)
% =========================================================
\begin{frame}\frametitle{Step (c): Simulation Lemma (Per-episode bound)}

\small
Let $d_h^{\pi^k}(s,a)=\Pr(s_h=s,a_h=a\mid s_0,\pi^k,P^\star)$ be the stage-$h$ occupancy.

\vspace{0.10cm}
On the uniform confidence event, for every episode $k$,
\[
V^\star - V^{\pi^k}
\ \le\
2\sum_{h=0}^{H-1}\E_{(s_h,a_h)\sim d_h^{\pi^k}}\!\big[b_h^k(s_h,a_h)\big].
\]

\vspace{0.12cm}
\pause
\textbf{Plug in the LinUCB-style bonus:}
\[
b_h^k(s,a)=\beta\,\|\phi(s,a)\|_{(\Lambda_h^k)^{-1}}
\quad\Longrightarrow\quad
V^\star - V^{\pi^k}
\ \le\
2\beta\sum_{h=0}^{H-1}\E_{d_h^{\pi^k}}\!\Big[\|\phi(s_h,a_h)\|_{(\Lambda_h^k)^{-1}}\Big].
\]

\end{frame}


% =========================================================
% Elliptical potential (make the log-det step explicit)
% =========================================================
\begin{frame}\frametitle{Summation: Elliptical Potential (Log-det)}

\small
Summing the simulation bound over episodes reduces regret to
\[
R_K \ \le\ 2\beta\sum_{h=0}^{H-1}\sum_{k=0}^{K-1}\|\phi(s_h^k,a_h^k)\|_{(\Lambda_h^k)^{-1}}.
\]

\vspace{0.10cm}
\pause
Fix a stage $h$ and write $x_k:=\phi(s_h^k,a_h^k)$, $\Lambda_{k+1}=\Lambda_k+x_kx_k^\top$.
Matrix determinant lemma gives the standard potential bound:
\[
\sum_{k=0}^{K-1}\|x_k\|_{\Lambda_k^{-1}}^2
\ \le\ 2\log\frac{\det(\Lambda_K)}{\det(\Lambda_0)}
\ \le\ 2d\log\!\Big(1+\tfrac{K}{\lambda}\Big).
\]

\vspace{0.10cm}
\pause
By Cauchy--Schwarz,
\[
\sum_{k=0}^{K-1}\|x_k\|_{\Lambda_k^{-1}}
\ \le\
\sqrt{K\sum_{k=0}^{K-1}\|x_k\|_{\Lambda_k^{-1}}^2}
\ \le\
\sqrt{2K\,d\,\log\!\Big(1+\tfrac{K}{\lambda}\Big)}.
\]
Sum over $h$ and multiply by $\beta$ to get the theorem.

\end{frame}

\section{Bridge: Bellman Rank \& Generalization}

\begin{frame}\frametitle{The Question: What makes RL Generalize?}

We have seen two successful "recipes" for avoiding $|\Scal|$-dependence:
\begin{enumerate}
    \item \textbf{Tabular UCBVI:} Estimate the full local transition.
    \item \textbf{Lin-UCBVI:} Estimate only the Bellman-relevant lookahead functionals.
\end{enumerate}

\vspace{0.4cm}
\pause
\textbf{The General Problem:}
What if we don't know the features $\phi(s,a)$? Or what if the transitions are non-linear, but some other structural bottleneck exists?

\vspace{0.4cm}
\pause
\textbf{The Goal:} We seek a \textbf{structural complexity measure} that tells us when trajectory data collected under policy $f$ can effectively "test" the quality of another hypothesis $g$.
\end{frame}


\begin{frame}\frametitle{The Core Object: Average Bellman Error}

Define the one-step Bellman residual of a hypothesis $g \in \mathcal{H}$ as:
\[
\ell_h(s,a,s';g) \;:=\; Q_{h,g}(s,a) - r_h(s,a) - V_{h+1,g}(s')
\]

\pause
\vspace{0.2cm}
We measure the quality of $g$ using its \textbf{Average Bellman Error} under the roll-in distribution induced by the policy of $f$, denoted $d_h^{\pi_f}$:
\[
\mathcal{E}_h^{Q}(f,g) \;:=\; \E_{(s,a)\sim d_h^{\pi_f},\ s'\sim P_h^\star}\Big[\ell_h(s,a,s';g)\Big]
\]

\pause
\vspace{0.2cm}
\textbf{Key Intuition:}
\begin{itemize}
    \item If $g = f^\star$ (realizability), then $\mathcal{E}_h(f, f^\star) = 0$ for \emph{every} roll-in policy $f$.
    \item Finding a near-optimal policy is equivalent to finding a $g$ that has low Bellman error under its own distribution ($\mathcal{E}_h(g, g) \approx 0$).
\end{itemize}
\end{frame}


\begin{frame}\frametitle{Defining Bellman Rank}

Consider a massive matrix $M_h$ where rows are "Roll-in Policies" $f \in \mathcal{H}$ and columns are "Tested Hypotheses" $g \in \mathcal{H}$. Each entry is the cross-error $\mathcal{E}_h(f,g)$.

\vspace{0.2cm}
\begin{definition}[$Q$-Bellman Rank]
The pair $(\Mcal,\Hcal)$ has \emph{$Q$-Bellman rank at most $d$} if there exist maps $X_h, W_h: \Hcal \to \R^d$ such that:
\[
\mathcal{E}_h^{Q}(f,g) \;=\; \langle X_h(f),\, W_h(g)\rangle \qquad \forall f,g\in\Hcal
\]
\end{definition}

\pause
\vspace{0.2cm}
\textbf{Conceptual Punchline:}
\begin{itemize}
    \item A Linear MDP assumes the \textbf{Transition Matrix} is low rank ($rank \le d$).
    \item Bellman Rank assumes the \textbf{Error Matrix} is low rank ($rank \le d$).
    \item We do not need to know the features $X_h, W_h$; we only assume they exist.
\end{itemize}
\end{frame}


\begin{frame}\frametitle{Example: LinMDPs have Bellman Rank $d$}

Why did today's Linear MDP model work? Let's check its Bellman Rank:
\begin{itemize}
    \item \textbf{LinMDP Hypothesis:} $Q_{h,g}(s,a) = \phi(s,a)^\top w_{h,g}$.
    \item \textbf{Bellman Error:} $\mathcal{E}_h^{Q}(f,g) = \E_{d_h^{\pi_f}}[ \phi(s,a)^\top w_{h,g} - r_h(s,a) - \E_{P^\star}[V_{h+1,g}(s')] ]$.
\end{itemize}

\pause
\vspace{0.2cm}
From our "Magic Property" Lemma: $r_h(s,a) + \E_{P^\star}[V_{h+1,g}]$ is linear in $\phi(s,a)$. Let its coefficient be $w_{lookahead}^\star$.
\begin{align*}
\mathcal{E}_h^{Q}(f,g) &= \E_{(s,a)\sim d_h^{\pi_f}} \Big[ \phi(s,a)^\top (w_{h,g} - w_{lookahead}^\star) \Big] \\
&= \Big\langle \underbrace{\E_{d_h^{\pi_f}}[\phi(s,a)]}_{X_h(f)}, \underbrace{w_{h,g} - w_{lookahead}^\star}_{W_h(g)} \Big\rangle
\end{align*}

\pause
\alert{Result:} LinMDPs are just one specific instance of the "Model Zoo" unified by Bellman Rank.
\end{frame}


\begin{frame}\frametitle{Algorithmic Teaser: Optimistic Elimination}

How do we learn if the error matrix is low-rank but the features are unknown?

\vspace{0.4cm}
\textbf{The PAC-RL Strategy:}
\begin{enumerate}
    \item \textbf{Eliminate:} If a hypothesis $g$ shows a large empirical Bellman error on any past roll-in data, throw it out.
    \item \textbf{Optimism:} Pick the $f_{t}$ that predicts the highest value among survivors.
    \item \textbf{Roll-in:} Collect new data using policy $\pi_{f_t}$.
\end{enumerate}

\vspace{0.4cm}
\pause
\textbf{The Logic for Next Time:}
Since the error matrix has rank $d$, we can only "fail" to eliminate a sub-optimal $f_t$ a maximum of $d$ times before we have spanned the space of possible errors. 

\vspace{0.2cm}
\textit{Next Lecture: Formalizing the "Simplicity" of RL and the PAC proof.}
\end{frame}

\end{document}

